% Elementos pré-textuais do Relatório de Estágio Supervisionado.
% Ordem e seções adaptadas do modelo Word disponibilizado pela UFAC.
% Os campos entre colchetes devem ser preenchidos antes da entrega.

% ---------------------------------------------------------
% Termo de aprovação (dados e assinaturas ainda pendentes).
% ---------------------------------------------------------
\begin{folhadeaprovacao}
  \begin{center}
    {\ABNTEXchapterfont\bfseries\large TERMO DE APROVAÇÃO\par}
    \vspace*{1.2cm}

    {\large\MakeUppercase{\imprimirautor}\par}
    \vspace*{1.2cm}

    {\ABNTEXchapterfont\bfseries\large\MakeUppercase{\imprimirtitulo}\par}
  \end{center}

  \vspace{1.3cm}
  Este relatório de estágio supervisionado foi apresentado ao Curso de
  Bacharelado em Sistemas de Informação da Universidade Federal do Acre
  (UFAC), como exigência parcial para obtenção do grau de bacharel
  em Sistemas de Informação.

  \vspace{0.6cm}
  \noindent\textbf{Data de aprovação:} [DIA/MÊS/ANO]

  \vspace{0.5cm}
  \noindent\textbf{Banca examinadora:}

  \assinatura{\textbf{Laura Costa Sarkis}\\Orientadora\\Universidade Federal do Acre}

  \assinatura{\textbf{[NOME COMPLETO DO(A) MEMBRO DA BANCA]}\\
    {[TITULAÇÃO E VÍNCULO INSTITUCIONAL]}}

  \begin{center}
    \vfill
    Rio Branco -- AC, [DIA] de [MÊS] de [ANO].
  \end{center}
\end{folhadeaprovacao}
\cleardoublepage

% ---------------------------------------------------------
% Dedicatória (opcional no texto final; espaço preservado).
% ---------------------------------------------------------
\begin{dedicatoria}
  \vspace*{\fill}
  \begin{flushright}
    \begin{minipage}{0.65\textwidth}
      \raggedleft
      \textit{[Inserir dedicatória, caso desejado.]}
    \end{minipage}
  \end{flushright}
  \vspace*{\fill}
\end{dedicatoria}
\cleardoublepage

% ---------------------------------------------------------
% Agradecimentos.
% ---------------------------------------------------------
\begin{agradecimentos}
  \noindent\textit{[Inserir os agradecimentos pessoais, acadêmicos e
  institucionais, conforme desejado.]}
\end{agradecimentos}
\cleardoublepage

% ---------------------------------------------------------
% Epígrafe (opcional no texto final; sem citação inventada).
% ---------------------------------------------------------
\begin{epigrafe}
  \vspace*{\fill}
  \begin{flushright}
    \begin{minipage}{0.65\textwidth}
      \raggedleft
      \textit{[Inserir frase ou citação escolhida.]}\\[0.6em]
      [AUTORIA DA EPÍGRAFE]
    \end{minipage}
  \end{flushright}
  \vspace*{\fill}
\end{epigrafe}
\cleardoublepage

% ---------------------------------------------------------
% Resumo e palavras-chave (a redigir após revisão dos resultados).
% ---------------------------------------------------------
\begin{resumo}
  \noindent
  \textit{[Redigir aqui um único parágrafo com a contextualização do
  estágio, o objetivo, a metodologia, as atividades desenvolvidas,
  os principais resultados comprovados e a conclusão.]}

  \vspace{\onelineskip}
  \noindent\textbf{Palavras-chave:} [Palavra-chave 1].
  [Palavra-chave 2]. [Palavra-chave 3].
\end{resumo}
\cleardoublepage

% ---------------------------------------------------------
% Abstract e keywords (tradução do resumo final).
% ---------------------------------------------------------
\begin{resumo}[Abstract]
  \begin{otherlanguage*}{english}
    \noindent
    \textit{[Write the English translation of the final Portuguese
    abstract here, including context, objectives, methods, findings,
    and conclusions.]}

    \vspace{\onelineskip}
    \noindent\textbf{Keywords:} [Keyword 1]. [Keyword 2]. [Keyword 3].
  \end{otherlanguage*}
\end{resumo}
\cleardoublepage

% ---------------------------------------------------------
% Listas automáticas, mesmo que ainda não existam ilustrações.
% Serão preenchidas pelo LaTeX conforme figuras, quadros e tabelas
% com legendas forem inseridos nos capítulos.
% ---------------------------------------------------------
\pdfbookmark[0]{\listfigurename}{lof}
\listoffigures*
\cleardoublepage

\pdfbookmark[0]{\listofquadrosname}{loq}
\listofquadros
\cleardoublepage

\pdfbookmark[0]{\listtablename}{lot}
\listoftables*
\cleardoublepage

% Siglas/abreviaturas existentes no texto. Revisar ao finalizar.
\begin{siglas}
  \item[API] Interface de Programação de Aplicações
  \item[IHC] Interação Humano-Computador
  \item[MVP] Produto Mínimo Viável
  \item[UFAC] Universidade Federal do Acre
\end{siglas}
\cleardoublepage

% Se o relatório passar a utilizar notação específica, inserir
% aqui um ambiente "simbolos" antes do sumário.
